
\((1+X^{(j)}t)_{j=1,\ldots,k}\) de \(\bigl(\Ghat(K[(X^{(j)})_j])\bigr)^k\)~; un
élément \(\theta\) de \(H(K[(X^{(j)})_j])\) détermine un morphisme de foncteurs si, et seulement si, il vérifie :
\[
\begin{array}{ll}
\text{i)}& \forall j,\quad H(f_1^{(j)})(\theta)=0\\[0.3em]
\text{ii)}& \forall n,j,\quad
H(f_n^{(j)})\bigl(u_H^{n,(j)}(\theta)\bigr)
\equiv u_H^{n,(j)}(\theta)\ \mathrm{mod.}\ F^nH.
\end{array}
\]

Le lecteur se convaincra aisément qu'à des complications de notation près, il suffit de recopier la démonstration de 1.6.

1.8. On notera \(\mathcal C_0\) la catégorie des \(K\)-algèbres, et
\[
  i:\mathcal C_0\longrightarrow \mathcal C,
\]
le foncteur qui à une \(K\)-algèbre \(A\) associe la \(K\)-algèbre \(A\) munie de la filtration grossière~; on notera
\[
  g:\mathcal C\longrightarrow\mathcal C_0,
\]
le foncteur «~oubli de la filtration~». Ces foncteurs vérifient les relations suivantes :
\[
  g\circ i=\operatorname{Id}_{\mathcal C_0};\qquad
  \Ghat\circ i\circ g=\Gg;\qquad
  \Gg\circ i\circ g=\Gg
\]
(cf. 1.3, ex. 2)

\textbf{Lemme 1.9.} \uline{Soit} \(H:\mathcal C_0\to\underline{\mathrm{Ab}}\,\underline{\mathrm{fil}}\) \uline{un foncteur sur} \(\mathcal C_0\)~; \uline{alors l'homomorphisme canonique}~:
\[
  \Hom_{\underline{\mathrm{Ab}}\,\underline{\mathrm{fil}}}(\Ghat,H\circ g)
  \longrightarrow
  \Hom_{\underline{\mathrm{Ab}}\,\underline{\mathrm{fil}}}(\Ghat\circ i,H)
\]
